% Generated draft: requires mathematical and source review.
\documentclass[11pt]{article}
\usepackage[margin=25mm]{geometry}
\usepackage{fontspec}
\setmainfont{Latin Modern Roman}
\usepackage{amsmath,amssymb,mathtools,unicode-math}
\setmathfont{Latin Modern Math}
\usepackage{enumitem,needspace,xurl,hyperref}
\setlength{\parindent}{0pt}
\setlength{\parskip}{5pt}
\setlength{\emergencystretch}{3em}
\newcommand{\N}{\mathbb{N}}
\newcommand{\Z}{\mathbb{Z}}
\newcommand{\Q}{\mathbb{Q}}
\newcommand{\R}{\mathbb{R}}
\newcommand{\C}{\mathbb{C}}
\newcommand{\F}{\mathbb{F}}
\newcommand{\A}{\mathbb{A}}
\newcommand{\PP}{\mathbb{P}}
\newcommand{\E}{\mathbb{E}}
\newcommand{\eps}{\varepsilon}
\newcommand{\dd}{\,\mathrm{d}}
\begin{document}
\section*{UnsolvedMath ID: 3100013}
\textbf{Source problem code:} catalogue entry
\subsection{Definitions and mathematical statement}
For a finite hypergraph $H=(V,E)$, its two-color discrepancy is
\[
\operatorname{disc}(H)=
\min_{\chi:V\to\{-1,1\}}\ \max_{e\in E}
\left|\sum_{v\in e}\chi(v)\right|.
\]
Set $\operatorname{disc}(H)=0$ when $E=\varnothing$. The source's literal assertion, attributed to Beck, is
\[
\exists c<\infty\ \forall\text{ finite hypergraphs }H:\qquad
\operatorname{disc}(H)\leq c\sqrt{|E(H)|}.
\]
The constant is universal, and the parameter is the \emph{number of hyperedges}, not their sizes, the number of vertices, or maximum vertex degree.

The supplied research summary reports this number-of-sets bound as a consequence of Spencer's theorem (1985), including sets of arbitrary finite size. It also cites the distinct Beck--Fiala theorem (1981): if each vertex lies in at most $t$ hyperedges, discrepancy is at most $2t-1$. That latter bound is linear in $t$, not a square-root-in-$t$ conclusion, and it is not replaced here by a stronger unprovided theorem.

The summary calls the list statement garbled, speculates about number-of-vertices or bounded-degree readings, and alternates between solved and PARTIAL-PROGRESS language even though the structured registry labels are solved and SOLVED-IN-LITERATURE. No additional hypotheses defining those alternative readings are supplied. In particular, a vertex-count bound for arbitrary families and a degree-based bound are different mathematical targets and cannot be identified with the literal displayed edge-count assertion. This draft preserves the literal inequality and the supplied claimed Spencer resolution, while leaving any different intended Beck problem unresolved.
\subsection{Short English statement}
The literal question bounds discrepancy by a universal constant times the square root of the number of hyperedges. The supplied summary reports that reading solved by Spencer, but calls the source garbled and speculates about different parameters. Those alternatives and the separate Beck--Fiala bound are kept distinct.
\subsection{Sources}
\begin{itemize}
\item Local UnsolvedMath registry, permanent ID 3100013, exact source code catalogue entry, ``Show that the discrepancy of any hypergraph H is at most c|E(H)|\^{}(1/2)''; proposer Beck; structured status solved, classification SOLVED-IN-LITERATURE. The supplied narrative also uses PARTIAL-PROGRESS for an unspecified corrected reading.
\item Cooper, \emph{Combinatorial Problems I Like}, October 2020 snapshot, section ``Graphs, Hypergraphs, Set Systems, Designs,'' source-order bullet 13, \url{https://people.math.sc.edu/cooper/combprob.html}, HTML extraction accessed July 29, 2026. The provenance says the author-maintained snapshot labels it open and later status was not checked item by item.
\item Spencer (1985), cited in the supplied summary for the square-root number-of-sets discrepancy bound; Beck--Fiala (1981), cited for the maximum-degree bound $2t-1$. No titles, complete bibliographic entries, or exact theorem statements are supplied.
\end{itemize}
\end{document}
